\ficha{Almeida (2010), Política econômica e crises cambiais}{almeida2010}

\begin{identificacao}
ALMEIDA, José Tadeu de. Política econômica e crises cambiais: a gestão financeira do Império do Brasil nos primórdios do padrão-ouro (1846-1858). \emph{História e Economia Revista Interdisciplinar}, 2010, p. 49-69.

Artigo, lido na íntegra, 22 páginas, numeração impressa de 49 a 69, bibliografia a partir da p. 67. Português. O PDF é separata sem folha de rosto: volume e número não constam.
\end{identificacao}

\bloco{Problema e tese}
Como as autoridades do Império conduziram a política monetária e cambial entre 1846 e 1858, e por que a defesa da paridade prevaleceu sobre a expansão do meio circulante. A resposta: a gestão conservadora subordinou a política econômica ao equilíbrio orçamentário e à valorização do mil-réis, tratando a emissão como recurso a evitar, porque a legitimidade internacional da paridade era condição da inserção periférica no padrão ouro-libra. Almeida acrescenta que essa ortodoxia servia à construção do Estado nacional, não à preservação de estruturas arcaicas (p. 51).

\chave{desenvolveu-se em sua maior parte sob a égide dos conservadores, cujo foco incidia constantemente sobre o equilíbrio orçamentário e pela manutenção da valorização da taxa de câmbio, como forma de evitar o recurso da emissão de moeda para sanear os gastos públicos}{p. 49}

\bloco{Argumento}
\begin{enumerate}[nosep]
\item O ingresso de 1846 responde a debilidade prévia, com café ainda sem surtos exportadores, moeda falsa e evasão de metálico (p. 54). A Lei nº 401 fixa 27 pence por mil-réis e autoriza o governo a recolher papel-moeda para sustentar a valorização (p. 55).
\item O padrão-ouro é lido como instituição de época, na chave de Eichengreen: prioridade à conversibilidade e à reserva que a sustente (p. 50 e p. 53). No plano internacional ele salvaguarda os grandes credores, notadamente as casas bancárias inglesas, e é esse interesse que condiciona políticas de cunho ortodoxo (p. 52). A assimetria é explícita: a Inglaterra manteve o padrão por quase um século, e as economias periféricas eram rotineiramente obrigadas a rever ou suspender suas metas de paridade-ouro (p. 56).
\item Daí a contradição posta como hipótese paralela: a expansão dos negócios internos pedia crédito e emissão, a legitimidade externa pedia defesa do mil-réis e restrição à livre emissão (p. 51). Os anos 1850 multiplicam casas bancárias de caráter emissionista (p. 56), e em 1857 a razão entre crédito e reservas chega a 300\% no Brasil, contra 100\% autorizados na França (p. 57).
\item A crise de 1857 mede a vulnerabilidade. Credores pressionam o Tesouro e o Banco do Brasil, o ouro embarca, e o governo suspende a conversibilidade entre notas do Tesouro e moeda metálica em 11 de novembro (p. 61). O socorro é empréstimo externo com aval do Império junto a Rothschild \& Sons, 600 mil libras, descrito como emprestador de última instância (p. 62). O câmbio cai a 23,5 pence em 31 de dezembro (p. 62) e só volta a 26,55 pence em agosto de 1858 (p. 63).
\item A correção é ortodoxa e rápida. Sousa Franco priorizara o meio circulante e a pluralidade de emissão (p. 59). Torres Homem, ministro em dezembro de 1858, revoga a autorização de emitir ao triplo do fundo disponível e retira a livre emissão dos bancos privados (p. 64). De 1858-59 em diante, as restrições progressivas marcam ponto de inflexão do Segundo Reinado (p. 56).
\end{enumerate}

\bloco{Conceitos e definições operacionais}
Padrão-ouro: sistema de paridades gerido internacionalmente, cuja pedra fundamental é a prioridade dada à conversibilidade e cuja função é manter ativa a solvência dos devedores frente aos credores (p. 52 e p. 53). Papel-moeda e moeda-papel: a nota 10 da p. 55 separa moeda emitida com lastro inferior a 100\% de moeda com lastro integral, e só a segunda seria compatível com a plena conversibilidade. Inserção periférica: adesão de jure à paridade sem capacidade de sustentá-la (p. 56 e p. 65). Valorização cambial: recolher papel-moeda e restringir a emissão (p. 55 e p. 64).

\chave{a partir da inserção de natureza periférica ao modelo de paridade cambial corporificado no padrão ouro-libra, então vigente, elucidando também a vulnerabilidade do sistema monetário brasileiro no século XIX}{p. 49}

\bloco{Evidência e método}
Narrativa histórica com sete gráficos de meio circulante, emissões, depósitos, reservas, câmbio e comércio exterior, montados sobre Ónody, Peláez e Suzigan, Villela e Mont'Alegre. Fontes de época em segunda mão, entre elas o relatório do Banco do Brasil de 1858. Recorte de 1846 a 1858, com foco nos gabinetes da Conciliação. Sem contrafactual: as séries ilustram o argumento, não o testam.

\bloco{Agentes e vertentes}
Itaboraí conduz a unificação bancária de 1853 (p. 57). Carneiro Leão autoriza em 1854 a emissão ao triplo do fundo disponível, para evitar a queda do câmbio (p. 58). Sousa Franco, do Partido Liberal, defende a pluralidade de emissão (p. 59). Salles Torres Homem tem concepção que o autor chama de corte nitidamente metalista (p. 63). Aparecem ainda Mauá, o Banco do Brasil e Rothschild \& Sons. Almeida recusa classificar os ministros como metalistas ou papelistas e propõe lê-los pelo interesse na construção centralizada do Estado (p. 66).

\bloco{Críticas e limites}
Texto de doutorando derivado da dissertação de mestrado do autor, e com marcas disso: a conclusão chama o próprio artigo de dissertação (p. 66). O aparato tem inconsistências, com Eichengreen citado ora como 2000 ora como 2002 e outras divergências entre corpo e bibliografia, o que obriga a conferir qualquer dado antes de reutilizá-lo. O artigo trata a Lei nº 401 como ingresso efetivo no regime do padrão-ouro (p. 54) sem discutir a literatura que nega vigência plena da conversibilidade no Império. A tese sobre a construção do Estado é afirmada mais do que demonstrada. O texto não trata da Primeira República nem da Caixa de Conversão.

\begin{dialogo}
\item Procurar em O Paiz e no Correio da Manhã, em 1906, invocações da Lei de 1846, dos 27 pence como par legal, e de Torres Homem ou Sousa Franco como antepassados do debate. Testa se os agentes de 1906 se viam como herdeiros de uma ortodoxia imperial ou se a filiação é construção da historiografia.
\item Procurar nos editoriais a ligação entre o limite de emissão da Caixa e a confiança dos credores de Londres. Testa a hipótese de que a ortodoxia se explica por necessidade de legitimidade internacional, e não por adesão doutrinária.
\item Procurar, nos debates sobre a custódia do ouro em 1911 e na suspensão de 1914, o arranjo de 1857, em que o Estado avaliza empréstimo externo ao banco emissor.
\item Verificar se 11 de novembro de 1857 é citado como precedente quando o troco da Caixa é suspenso em 1914. Testa se a suspensão era lida como falência da regra ou como cláusula prevista.
\end{dialogo}

\usonocapitulo{Parte III, como apoio. Mostra que a ortodoxia de 1898 a 1906 tinha precedente imperial: equilíbrio orçamentário, valorização do mil-réis e restrição da emissão já formavam pacote coerente sob os gabinetes conservadores, justificado pela inserção periférica no padrão ouro-libra e pelo interesse dos credores ingleses. Serve também ao argumento de vulnerabilidade, já que a suspensão de 1857 mostra a regra como contingente na periferia desde o início. Não serve como evidência sobre 1906, que o texto não aborda, nem como autoridade em dados quantitativos.}
